\chapter{अवकल रूप और स्टोक्स की प्रमेय}\label{ch:b3:forms}

विश्लेषण की एक प्रमेय ने दो शताब्दियों में अपनी जाति की शेष सभी
प्रमेयों को अपने में समा लिया है: कलन की मूल प्रमेय, ग्रीन--रीमान
(दूसरे~वर्ष के खंड में सिद्ध), गाउस की अपसरण प्रमेय,
केल्विन--स्टोक्स की कर्ल प्रमेय --- इनमें से प्रत्येक कहती है कि
किसी क्षेत्र पर किसी अवकलज का समाकल \omterm{def:b3:forms:boundary}{परिसीमा} पर मूल वस्तु के समाकल
के बराबर है। \emph{\omterm{def:b3:forms:diffform}{अवकल रूपों}} की भाषा उन सबको एक ही कथन $\int_M\dd\omega = \int_{\partial M}\omega$ बना
देती है, और उस कथन को एक ही आघात में सिद्ध करने योग्य भी। यह अध्याय
वह भाषा ईमानदारी से बनाता है --- एकांतर बहुरैखिक बीजगणित, \omterm{def:b3:forms:d}{बाह्य
अवकलज}, \omterm{def:b3:forms:pullback}{प्रत्यानयन}, \omterm{def:b3:forms:orientation}{अभिविन्यास}, \cref{ch:b3:submanifolds} के उपबहुविधों पर समाकलन ---
स्टोक्स की प्रमेय सिद्ध करता है, और पहले चेक भुनाता है: शास्त्रीय
समाकल प्रमेय, वह \omterm{def:b3:forms:winding}{घूर्णन संख्या} जो चुपचाप \cref{ch:b3:residues} चला रही थी, और
सप्ताहांत समस्या में ब्राउवर की स्थिर-बिंदु प्रमेय। सर्वत्र
\emph{चिकना} से अभिप्राय है $\mathcal
C^\infty$; और जब तक अन्यथा न कहा जाए,
प्रत्येक प्रतिचित्रण और रूप चिकना है। इससे इस स्तर पर काम की कोई
व्यापकता नहीं खोती और हाथ खुल जाते हैं।

\section{एकांतर बहुरैखिक बीजगणित}

\begin{definition}\label{def:b3:forms:alternating}
मान लीजिए $E$ विमा $n$ की कोई वास्तविक सदिश समष्टि है। $E$ पर
\emph{$k$-रैखिक एकांतर रूप}\index{एकांतर रूप} ऐसा प्रतिचित्रण
$\alpha\colon E^k \to \R$ है जो प्रत्येक चर में रैखिक हो और जिसके लिए दो कोटियाँ बराबर
होने पर $\alpha(v_1, \dots, v_k) = 0$ हो। उनकी समष्टि $\Lambda^k E^*$ लिखी जाती है; और परिपाटी से
$\Lambda^0E^* = \R$। एकांतरता प्रतिसममिति अनिवार्य कर देती है: दो कोटियों की
अदला-बदली चिह्न बदल देती है ($\alpha(\dots, v + w, \dots, v + w, \dots) = 0$ का प्रसार कीजिए), और अधिक व्यापक
रूप से प्रत्येक क्रमचय $\sigma$ के लिए $\alpha(v_{\sigma(1)}, \dots, v_{\sigma(k)}) =
\varepsilon(\sigma)\,\alpha(v_1, \dots, v_k)$।
\end{definition}

\begin{example}
$E = \R^n$ पर: $\Lambda^1E^* = E^*$ \omterm{def:b3:banach:operator}{द्वैत समष्टि} है; विहित आधार में सारणिक एक
$n$-रैखिक \omterm{def:b3:forms:alternating}{एकांतर रूप} है, और \cref{prop:b3:forms:basis} दिखा देगा कि वह $\Lambda^nE^*$ को फैला
देता है --- यही वह गहरा कारण है कि सारणिक मापन तक अद्वितीय है।
$k > n$ के लिए $\Lambda^kE^* = \{0\}$: $k$ सदिश आश्रित होते हैं, और एक को दूसरों के
अनुदिश फैलाने पर एकांतरता $\alpha$ को मार देती है।
\end{example}

\begin{definition}\label{def:b3:forms:wedge}
$\ell_1, \dots, \ell_k \in E^*$ के लिए उनका \emph{बाह्य गुणन}\index{बाह्य गुणन} $k$-रैखिक
\omterm{def:b3:forms:alternating}{एकांतर रूप}
\[
(\ell_1 \wedge \dots \wedge \ell_k)(v_1, \dots, v_k)
= \det\bigl(\ell_i(v_j)\bigr)_{1 \leq i, j \leq k} .
\]
है। एकांतरता और बहुरैखिकता सारणिक की, उसके स्तंभों में, वही हैं।
\end{definition}

\begin{proposition}[$\Lambda^k$ का आधार]\label{prop:b3:forms:basis}
मान लीजिए $(e_1, \dots, e_n)$ $E$ का कोई आधार है जिसका द्वैत आधार $(e_1^*, \dots, e_n^*)$ है। तब
रूप
\[
e_I^* = e_{i_1}^* \wedge \dots \wedge e_{i_k}^*,
\qquad I = \{i_1 < \dots < i_k\} \subseteq \{1, \dots, n\},
\]
$\Lambda^kE^*$ का आधार बनाते हैं; अतः $\dim\Lambda^kE^* =
\binom nk$। स्पष्ट रूप से, $\alpha =
\sum_{\abs I = k}\alpha(e_{i_1}, \dots, e_{i_k})\,e_I^*$।
\end{proposition}

\begin{proof}
\emph{जनकता।} मान लीजिए $\alpha \in \Lambda^kE^*$ और $\beta =
\sum_I\alpha(e_I)\,e_I^*$, जहाँ $\alpha(e_I)$ $\alpha(e_{i_1}, \dots, e_{i_k})$ का संक्षेप
है। दोनों पक्ष $k$-रैखिक और एकांतर हैं, अतः वे तभी सहमत हो जाते
हैं जब $j_1 < \dots <
j_k$ वाले सभी $k$-बहुकों $(e_{j_1}, \dots, e_{j_k})$ पर सहमत हों (बहुरैखिकता
आधार सदिशों के बहुकों तक सिमटा देती है, और एकांतरता कड़ाई से बढ़ते
हुए बहुकों तक)। और $e_I^*(e_{j_1},
\dots, e_{j_k}) = \det(e_{i_r}^*(e_{j_s})) = \delta_{IJ}$: $I = J$ के लिए आव्यूह तत्समक है; और $I \neq J$ के
लिए कोई पंक्ति शून्य होती है। अतः सभी $J$ के लिए $\beta(e_J) = \alpha(e_J)$: इसलिए
$\beta = \alpha$। \emph{स्वतंत्रता।} यदि $\sum_I c_Ie_I^* = 0$ हो, तो $(e_{j_1}, \dots, e_{j_k})$ पर मान लेने से
$c_J = 0$ मिलता है।
\end{proof}

\begin{definition}\label{def:b3:forms:wedgegeneral}
\omterm{def:b3:forms:wedge}{बाह्य गुणन} द्विरैखिक प्रतिचित्रण $\Lambda^kE^* \times \Lambda^\ell E^* \to
\Lambda^{k+\ell}E^*$ तक बढ़ जाता है, जो
द्विरैखिकता और $(e_I^*) \wedge (e_J^*) = e_I^* \wedge e_J^*$ से निर्धारित होता है (जोड़िए और पुनःक्रमित
कीजिए; $I \cap J \neq
\varnothing$ होने पर गुणनफल $0$ है)। वह साहचर्य और
\emph{श्रेणीबद्ध-प्रतिक्रमविनिमेय} है:
\[
\beta \wedge \alpha = (-1)^{k\ell}\,\alpha \wedge \beta
\qquad (\alpha \in \Lambda^k, \ \beta \in \Lambda^\ell).
\]
\end{definition}

\begin{proof}[\omnameProof]
दोनों गुण आधार अवयवों पर जाँचे जाते हैं और द्विरैखिकता से बढ़ा
दिए जाते हैं। साहचर्यता: \cref{def:b3:forms:wedge} के सारणिक सूत्र (खंडों द्वारा
लाप्लास प्रसार) से $e_I^* \wedge e_J^* \wedge e_K^*$ के दोनों कोष्ठकीकरण $1$-रूपों के जुड़े
हुए कुल के \omterm{def:b3:forms:wedge}{बाह्य गुणन} के बराबर होते हैं। चिह्न नियम: $\beta$ के
$\ell$ गुणनखंडों में से प्रत्येक को $\alpha$ के $k$ गुणनखंडों के पार
ले जाने पर प्रत्येक निकटवर्ती स्थानांतरण (सारणिक की दो पंक्तियों
की अदला-बदली) एक चिह्न की कीमत लेता है, अतः कुल $(-1)^{k\ell}$।
\end{proof}

\begin{proposition}[प्रत्यानयन, रैखिक स्थिति]
\label{prop:b3:forms:linearpullback}
कोई रैखिक प्रतिचित्रण $u\colon E \to F$ प्रत्येक $k$ के लिए रैखिक प्रतिचित्रण
$u^*\colon \Lambda^kF^* \to \Lambda^kE^*$, $(u^*\alpha)(v_1, \dots, v_k) = \alpha(u(v_1), \dots,
u(v_k))$, प्रेरित करता है। वह $u^*(\alpha \wedge \beta) = u^*\alpha
\wedge u^*\beta$ और $(u \circ w)^* = w^* \circ u^*$ संतुष्ट करता है।
इसके अतिरिक्त:
\begin{enumerate}
  \item यदि $\dim E = n$ और $u\colon E \to E$ हो, तो रेखा $\Lambda^nE^*$ पर: $u^*\alpha = (\det u)\,\alpha$।
  \item यदि $\operatorname{rk}u < k$ हो, तो $\Lambda^kF^*$ पर $u^* = 0$।
\end{enumerate}
\end{proposition}

\begin{proof}
कार्यकीय सर्वसमिकाएँ परिभाषाओं से तत्काल निकलती हैं (गुणनफल नियम
के लिए $1$-रूपों के बाह्य गुणनों पर सारणिक सूत्र से जाँचिए ---
$\det(\ell_i(uv_j)) =
\det((u^*\ell_i)(v_j))$ --- फिर द्विरैखिक रूप से बढ़ाइए)।
(1) $u^*$ एक-विमीय $\Lambda^nE^*$ (\cref{prop:b3:forms:basis}) को उसी पर भेजता है, अतः $u^*\alpha =
c\,\alpha$, जहाँ
$c$ $\alpha \neq 0$ से स्वतंत्र है; $\alpha = e_1^* \wedge \dots \wedge e_n^*$ और $(v_j) =
(e_j)$ पर परीक्षा लेने से $c = \det(e_i^*(ue_j)) = \det u$
मिलता है।
(2) $v_1, \dots, v_k \in E$ के लिए सदिश $u(v_1), \dots,
u(v_k)$ $u$ के प्रतिबिंब में होते हैं, जिसकी
विमा $< k$ है: अतः वे रैखिकतः आश्रित हैं, और कोई \omterm{def:b3:forms:alternating}{एकांतर रूप} किसी
आश्रित कुल पर लुप्त हो जाता है (आश्रित सदिश को दूसरों के अनुदिश
फैलाइए)।
\end{proof}

\section{अवकल रूप और बाह्य अवकलज}

\begin{definition}\label{def:b3:forms:diffform}
मान लीजिए $U \subseteq \R^n$ विवृत है। $U$ पर \emph{अवकल
$k$-रूप}\index{अवकल रूप} कोई चिकना प्रतिचित्रण $\omega\colon U \to \Lambda^k(\R^n)^*$ है; \cref{prop:b3:forms:basis}
के आधार में (जहाँ $e_i^*$ के लिए $\dd x_i$ लिखा जाता है),
\[
\omega = \sum_{\abs I = k} a_I\,\dd x_I,
\qquad
\dd x_I = \dd x_{i_1} \wedge \dots \wedge \dd x_{i_k},
\]
जिसमें गुणांक $a_I \in \mathcal C^\infty(U)$ चिकने हैं। उनकी समष्टि $\Omega^k(U)$ है; और $\Omega^0(U) = \mathcal
C^\infty(U)$।
कोई $0$-रूप एक फलन है; कोई $1$-रूप रैखिक रूपों का क्षेत्र है
(उदाहरणार्थ किसी फलन का अवकलज $\dd f$); और कोई $n$-रूप $a\,\dd x_1\wedge\dots\wedge\dd
x_n$ है,
जो \cref{ch:b3:product} का स्वाभाविक समाकल्य है।
\end{definition}

\begin{definition}[बाह्य अवकलज]\label{def:b3:forms:d}
\emph{बाह्य अवकलज}\index{बाह्य अवकलज} वह रैखिक प्रतिचित्रण $\dd\colon \Omega^k(U) \to \Omega^{k+1}(U)$
है जो
\[
\dd\Bigl(\sum_I a_I\,\dd x_I\Bigr)
= \sum_I \dd a_I \wedge \dd x_I
= \sum_I\sum_{j=1}^n \frac{\partial a_I}{\partial
x_j}\,\dd x_j \wedge \dd x_I .
\]
से परिभाषित होता है। $0$-रूपों पर वह सामान्य अवकलज ही है।
\end{definition}

\begin{theorem}\label{thm:b3:forms:dsquare}
(क) $\omega \in \Omega^k$ के लिए $\dd(\omega \wedge \eta) = \dd\omega \wedge \eta +
(-1)^k\,\omega \wedge \dd\eta$ (\emph{श्रेणीबद्ध लाइब्नित्स नियम})।
(ख) $\dd \circ \dd = 0$।\index{डी वर्ग शून्य@$\dd^2 = 0$}
\end{theorem}

\begin{proof}
(क) द्विरैखिकता से केवल $\omega = a\,\dd
x_I$, $\eta = b\,\dd x_J$ की स्थिति देखनी पर्याप्त है।
तब $\omega \wedge \eta =
ab\,\dd x_I \wedge \dd x_J$ और
\[
\dd(\omega\wedge\eta) = (b\,\dd a + a\,\dd b) \wedge \dd x_I
\wedge \dd x_J
= (\dd a \wedge \dd x_I) \wedge (b\,\dd x_J) +
a\,\dd b \wedge \dd x_I \wedge \dd x_J,
\]
और $\dd b$ $1$-रूप को $\dd
x_I$ के $k$ गुणनखंडों के पार ले जाने पर
$(-1)^k$ (\cref{def:b3:forms:wedgegeneral}) की कीमत लगती है: अतः दूसरा पद $(-1)^k\,\omega \wedge \dd\eta$ है।
(ख) किसी फलन के लिए: $\dd(\dd a) = \sum_{i,j}
\frac{\partial^2a}{\partial x_j\partial x_i}\,\dd x_j \wedge
\dd x_i$। श्वार्ट्स की मिश्रित आंशिक अवकलजों
संबंधी प्रमेय (दूसरे~वर्ष के खंड में सिद्ध; $a$ $\mathcal C^\infty$ है) से
गुणांक $(i,j)$ में सममित है, जबकि $\dd x_j \wedge
\dd x_i$ प्रतिसममित है: पदों $(i,j)$
और $(j,i)$ को युग्मित करने पर सब कुछ निरस्त हो जाता है। किसी व्यापक
$\omega = \sum
a_I\dd x_I$ के लिए: (क) से $\dd\dd\omega = \sum\dd(\dd a_I \wedge \dd x_I)
= \sum(\dd\dd a_I\wedge\dd x_I - \dd a_I\wedge\dd(\dd x_I))$, और दोनों पद लुप्त हो जाते हैं ($\dd(\dd x_I) = 0$,
क्योंकि गुणांक अचर है)।
\end{proof}

\begin{definition}[प्रत्यानयन]\label{def:b3:forms:pullback}
मान लीजिए $\varphi\colon U \to V$ चिकना है ($U \subseteq \R^m$, $V \subseteq \R^n$ विवृत)।
\emph{प्रत्यानयन}\index{प्रत्यानयन} $\varphi^*\colon \Omega^k(V) \to \Omega^k(U)$ बिंदुशः अवकलज के
अनुदिश रैखिक प्रत्यानयन से परिभाषित होता है: $(\varphi^*\omega)_x = (D\varphi(x))^*\,\omega_{\varphi(x)}$। मूर्त रूप
में, $\varphi^*$ प्रतिस्थापन कर देता है: फलनों पर $\varphi^*f = f \circ
\varphi$, $\varphi^*(\dd y_i) =
\dd\varphi_i = \sum_j\frac{\partial\varphi_i}{\partial
x_j}\dd x_j$, और
$\varphi^*(a\,\dd y_{i_1}\wedge\dots\wedge
\dd y_{i_k}) =
(a\circ\varphi)\,\dd\varphi_{i_1}\wedge\dots\wedge
\dd\varphi_{i_k}$।
\end{definition}

\begin{theorem}\label{thm:b3:forms:pullback}
(क) $\varphi^*(\omega \wedge \eta) = \varphi^*\omega \wedge
\varphi^*\eta$ और $(\psi \circ \varphi)^* = \varphi^* \circ
\psi^*$।
(ख) $\varphi^*(\dd\omega) = \dd(\varphi^*\omega)$: अर्थात् \omterm{def:b3:forms:d}{बाह्य अवकलज} प्रत्येक चिकने प्रतिस्थापन के साथ
क्रमविनिमेय है --- यही वह सर्वसमिका है जो उसे इस सिद्धांत का
\emph{वही} अवकलज बना देती है।
(ग) यदि $\varphi\colon U \to V$ $\R^n$ और $\omega = a\,\dd y_1 \wedge \dots \wedge \dd
y_n$ के विवृत समुच्चयों के बीच चिकना हो,
तो $\varphi^*\omega = (a \circ
\varphi)\,\det\bigl(D\varphi\bigr)\,\dd x_1 \wedge \dots
\wedge \dd x_n$।
\end{theorem}

\begin{proof}
(क) ये रैखिक प्रत्यानयनों के बारे में बिंदुशः कथन हैं (\cref{prop:b3:forms:linearpullback}), और
साथ में शृंखला नियम $D(\psi\circ\varphi)(x) = D\psi(\varphi(x))\,D\varphi(x)$।
(ख) किसी $0$-रूप $f$ के लिए: $\varphi^*(\dd f) = \dd f \circ
D\varphi = \dd(f \circ \varphi)$ शृंखला नियम ही है। $\omega = a\,\dd y_I$
के लिए: (क) से $\varphi^*\omega =
(a\circ\varphi)\,\dd\varphi_{i_1}\wedge\dots\wedge
\dd\varphi_{i_k}$, अतः लाइब्नित्स नियम (\cref{thm:b3:forms:dsquare}(क)) और $\dd\dd\varphi_{i_r} =
0$ से:
\[
\dd(\varphi^*\omega) = \dd(a\circ\varphi) \wedge
\dd\varphi_{i_1}\wedge\dots\wedge\dd\varphi_{i_k}
= \varphi^*(\dd a) \wedge \varphi^*(\dd y_I)
= \varphi^*(\dd a \wedge \dd y_I) =
\varphi^*(\dd\omega).
\]
(ग) बिंदुशः यह उच्चतम-घात रूपों पर ठीक $u^*\alpha = (\det u)\alpha$ है ($u = D\varphi(x)$ के साथ
\cref{prop:b3:forms:linearpullback}(1))।
\end{proof}

\begin{example}[ध्रुवीय निर्देशांक]\label{ex:b3:forms:polar}
$\varphi(r, \theta) = (r\cos\theta, r\sin\theta)$ के लिए: $\varphi^*\dd x = \cos\theta\,\dd r -
r\sin\theta\,\dd\theta$, $\varphi^*\dd y = \sin\theta\,\dd r
+ r\cos\theta\,\dd\theta$, अतः
\[
\varphi^*(\dd x \wedge \dd y) =
(\cos\theta\,\dd r - r\sin\theta\,\dd\theta) \wedge
(\sin\theta\,\dd r + r\cos\theta\,\dd\theta)
= r\,\dd r \wedge \dd\theta,
\]
अर्थात् \cref{ex:b3:product:polar} का याकोबीय विशुद्ध बीजगणित से प्रकट हो जाता है ---
कोई \omterm{def:b3:measure:measure}{माप} सिद्धांत नहीं। \cref{exo:b3:forms:9} इस टिप्पणी को एक कथन में बदल देती
है: अभिविन्यस्त समाकलों के लिए चर परिवर्तन सूत्र \emph{वस्तुतः}
\omterm{def:b3:forms:pullback}{प्रत्यानयन} सूत्र ही है।
\end{example}

\section{संवृत और यथार्थ रूप; पोंकारे प्रमेयिका}

\begin{definition}\label{def:b3:forms:closedexact}
$\omega \in \Omega^k(U)$ \emph{संवृत}\index{संवृत रूप} कहलाता है यदि $\dd\omega = 0$ हो, और
\emph{यथार्थ}\index{यथार्थ रूप} यदि किसी $\eta \in \Omega^{k-1}(U)$ के लिए $\omega = \dd\eta$ हो
($\omega$ का कोई \emph{आद्यंतर})। $\dd^2 = 0$ से यथार्थ $\Rightarrow$ संवृत; और
विलोम $U$ के \emph{आकार} के बारे में एक प्रश्न है।
\end{definition}

\begin{example}[कोणीय रूप]\label{ex:b3:forms:angular}
$U = \R^2 \setminus \{0\}$ पर
\[
\omega_\theta = \frac{x\,\dd y - y\,\dd x}{x^2 + y^2}
\]
\omterm{def:b3:forms:closedexact}{संवृत} है (सीधा परिकलन: \cref{exo:b3:forms:4}) पर \omterm{def:b3:forms:closedexact}{यथार्थ} नहीं: इकाई वृत्त के
अनुदिश उसका समाकल $2\pi \neq 0$ है, जबकि \omterm{def:b3:forms:closedexact}{संवृत} वक्रों के अनुदिश \omterm{def:b3:forms:closedexact}{यथार्थ}
रूपों के समाकल लुप्त हो जाते हैं (\cref{prop:b3:forms:exactloop})। स्थानीय रूप से ध्रुवीय
कोण के किसी भी चिकने निर्धारण $\theta$ के लिए $\omega_\theta = \dd\theta$ --- और यहीं से
नाम तथा बाधा दोनों आते हैं: समूचे $U$ पर ऐसा कोई निर्धारण
विद्यमान नहीं है। यही एक रूप \omterm{def:b3:forms:winding}{घूर्णन संख्या} (\cref{sec:b3:forms:winding}) को चलाता है और,
उसके माध्यम से, \cref{ch:b3:residues} की \omterm{def:b3:residues:singularities}{अवशेष} प्रमेय को भी।
\end{example}

\begin{theorem}[पोंकारे प्रमेयिका]\label{thm:b3:forms:poincare}
मान लीजिए $U \subseteq \R^n$ विवृत और $0$ के परितः \emph{तारकाकार} है।
$U$ पर प्रत्येक \omterm{def:b3:forms:closedexact}{संवृत} $k$-रूप ($k \geq 1$) \omterm{def:b3:forms:closedexact}{यथार्थ}
होता है।\index{पोंकारे प्रमेयिका}
\end{theorem}

\begin{proof}
हम एक रैखिक \emph{समस्थेय संकारक} $h\colon
\Omega^k(U) \to \Omega^{k-1}(U)$ बनाते हैं जिसके लिए
\begin{equation}\label{eq:b3:forms:homotopy}
\dd(h\omega) + h(\dd\omega) = \omega
\qquad (k \geq 1);
\end{equation}
हो; तब $\dd\omega = 0$ होने पर $\omega = \dd(h\omega)$, और काम पूरा। $\omega = \sum_I a_I\,\dd x_I$ के लिए
\[
h\omega = \sum_{\abs I = k}\ \sum_{r=1}^{k}(-1)^{r-1}
\Bigl(\int_0^1 t^{k-1}a_I(tx)\,\dd t\Bigr)\,
x_{i_r}\,\dd x_{i_1}\wedge\dots\wedge
\widehat{\dd x_{i_r}}\wedge\dots\wedge\dd x_{i_k}
\]
रखिए (टोपी किसी गुणनखंड को हटा देती है; समाकल $x$ में चिकने हैं,
समाकल के भीतर अवकलन से, \cref{thm:b3:lebesgue:paramdiff}, और सभी अवकलज संहतों पर प्रभावित
हैं)। \eqref{eq:b3:forms:homotopy} की जाँच जीवन में एक ही बार का परिकलन है, अतः हम उसे
पूरा करते हैं। $I$ स्थिर कीजिए और $\omega =
a\,\dd x_I$ लीजिए (रैखिकता)। पहले,
\[
\dd(h\omega) = k\Bigl(\int_0^1t^{k-1}a(tx)\dd t\Bigr)\dd x_I
+ \sum_{r=1}^k(-1)^{r-1}\sum_{j=1}^n
\Bigl(\int_0^1t^{k}\,\partial_ja(tx)\,\dd t\Bigr)
x_{i_r}\,\dd x_j\wedge\dd x_{I\setminus i_r} :
\]
जहाँ पहला समूह उन पदों को इकट्ठा करता है जिनमें $\dd$ गुणनखंड
$x_{i_r}$ पर पड़ता है --- \omterm{def:b3:forms:wedge}{बाह्य गुणन} $\dd x_{i_r}\wedge\dd
x_{I\setminus i_r}$ चिह्न $(-1)^{r-1}$ के साथ $\dd x_I$
को फिर से जोड़ देता है, जो पूर्व-गुणक को निरस्त कर देता है, और
$r$ के $k$ मान गुणक $k$ दे देते हैं --- जबकि दूसरा समूह उन
पदों को इकट्ठा करता है जहाँ $\dd$ समाकल पर पड़ता है (शृंखला नियम
$t\,\partial_ja(tx)$ बाहर निकाल देता है)। इसके बाद $\dd\omega =
\sum_j\partial_ja\,\dd x_j\wedge\dd x_I$, और घात $k+1$ में $h$
की परिभाषा लगाने पर, जिसमें सूचकांक $j$ पहला स्थान लेता है:
\[
h(\dd\omega) = \sum_{j=1}^n\Bigl(\int_0^1t^{k}
\partial_ja(tx)\dd t\Bigr)x_j\,\dd x_I
- \sum_{j=1}^n\sum_{r=1}^k(-1)^{r-1}
\Bigl(\int_0^1t^{k}\partial_ja(tx)\dd t\Bigr)
x_{i_r}\,\dd x_j\wedge\dd x_{I\setminus i_r} .
\]
$\dd(h\omega) + h(\dd\omega)$ में द्विक योग निरस्त हो जाते हैं, जो इसलिए कलन की मूल प्रमेय
से
\[
\Bigl(\int_0^1\bigl(k\,t^{k-1}a(tx) +
t^k{\textstyle\sum_j}x_j\,\partial_ja(tx)\bigr)\dd t\Bigr)
\dd x_I
= \Bigl(\int_0^1\frac{\dd}{\dd t}\bigl(t^ka(tx)\bigr)\dd
t\Bigr)\dd x_I
= a(x)\,\dd x_I = \omega,
\]
के बराबर है। तारकाकारता वहीं आई जहाँ आनी चाहिए थी: $t \in \intcc01$ के लिए
$tx \in U$, ताकि $a(tx)$ सार्थक हो।
\end{proof}

\begin{remark}
$k = 1$ और $\omega = \sum_j a_j\dd x_j$ के लिए आद्यंतर $f(x) = \int_0^1\sum_ja_j(tx)\,x_j\,\dd t$ है --- अर्थात् रेखाखंड $[0, x]$
के अनुदिश $\omega$ का रेखा समाकल: यह प्रमेय दूसरे~वर्ष के
``सममित याकोबीय वाला क्षेत्र प्रवणक होता है'' का बहु-चर रूप है,
अब प्रत्येक घात में। कोणीय रूप (\cref{ex:b3:forms:angular}) दिखा देता है कि $U$ पर
परिकल्पना सजावट नहीं है: $\R^2\setminus\{0\}$ तारकाकार नहीं है, और वहाँ संवृतता
से यथार्थता नहीं निकलती। किसी व्यापक \omterm{def:b3:topology:topology}{विवृत समुच्चय} पर जो बचता है
उसे \emph{दे राम सह-समजातता} $H^k(U) =
\ker\dd/\operatorname{im}\dd$ मापती है --- पहली अतुच्छ गणना
के लिए \cref{exo:b3:forms:12} देखिए।
\end{remark}

\section{अभिविन्यास और उपबहुविधों पर समाकलन}

किसी $k$-रूप के समाकलन के लिए $k$-विमीय अभिविन्यस्त भूमि
चाहिए। \cref{ch:b3:submanifolds} (\cref{thm:b3:submanifolds:characterizations}) से स्मरण कीजिए कि कोई $k$-उपबहुविध $M \subseteq \R^n$
स्थानीय रूप से किसी नियमित प्राचलन $\gamma\colon V \to M \cap W$ का प्रतिबिंब होता है
($V
\subseteq \R^k$ विवृत, $\gamma$ एकैकी अवकलज के साथ अपने प्रतिबिंब पर
\omterm{def:b3:topology:continuity}{समस्थितिकता})।

\begin{definition}\label{def:b3:forms:orientation}
$M$ का \emph{अभिविन्यास}\index{अभिविन्यास} प्रत्येक $p \in M$ के
लिए \omterm{def:b3:submanifolds:tangent}{स्पर्श समष्टि} $T_pM$ के आधारों के दो अभिविन्यास वर्गों में से
एक का ऐसा चयन है जो \emph{स्थानीय रूप से संगत} हो: प्रत्येक बिंदु
के आसपास कोई ऐसा प्राचलन $\gamma$ हो जिसकी निर्देशांक चौखट $(\partial_1\gamma, \dots,
\partial_k\gamma)$
अपने प्रांत के प्रत्येक बिंदु पर धनात्मक रूप से अभिविन्यस्त हो।
ऐसे प्राचलन \emph{प्रत्यक्ष} कहलाते हैं। $M$
\emph{अभिविन्यासनीय} कहलाती है यदि कोई अभिविन्यास विद्यमान हो;
मोबियस पट्टी दिखा देती है कि यह विफल भी हो सकता है। इस अध्याय की
सभी उपबहुविध अभिविन्यस्त हैं।
\end{definition}

\begin{definition}[किसी रूप का समाकल]\label{def:b3:forms:integral}
मान लीजिए $M$ कोई अभिविन्यस्त $k$-उपबहुविध है और $\omega$ $M$
के किसी \omterm{def:b3:topology:topology}{प्रतिवेश} पर परिभाषित कोई $k$-रूप, जिसमें $\operatorname{supp}\omega \cap M$ \omterm{def:b3:topology:compact}{संहत} हो।
(क) यदि किसी एक ही प्रत्यक्ष प्राचलन के लिए $\operatorname{supp}\omega \cap M \subseteq
\gamma(V)$ हो, तो
\[
\int_M\omega = \int_V\gamma^*\omega
\]
रखिए --- जहाँ दायाँ पक्ष $V$ (\cref{ch:b3:product}) पर $\gamma^*\omega = g\,\dd u_1\wedge\dots\wedge\dd u_k$ के गुणांक $g$
का लेबेग समाकल है, जो \omterm{def:b3:topology:compact}{संहत} आलंब के साथ \omterm{def:b3:topology:continuity}{संतत} है।
(ख) व्यापक स्थिति में \omterm{def:b3:topology:compact}{संहत} $\operatorname{supp}\omega \cap M$ को आच्छादित करने वाले परिमित
संख्या के प्रत्यक्ष प्राचलन $\gamma_i(V_i)$ और उनके अधीनस्थ \omterm{lem:b3:forms:partition}{इकाई विभाजन}
$(\chi_i)$ (\cref{lem:b3:forms:partition}) चुनिए, और $\int_M\omega =
\sum_i\int_M\chi_i\,\omega$ रखिए, जिसका प्रत्येक पद (क) से
परिकलित होता है। किसी वक्र ($k = 1$) के लिए, जो $\gamma\colon\intcc ab\to\R^n$ से प्राचलित
हो, हम $\int_\gamma\omega = \int_a^b\gamma^*\omega$ लिखते हैं, और इसके लिए किसी एकैकीपन की आवश्यकता नहीं।
\end{definition}

\begin{lemma}[संगति]\label{lem:b3:forms:consistency}
परिभाषा (क) प्रत्यक्ष प्राचलन पर निर्भर नहीं करती, और परिभाषा (ख)
न आवरण पर निर्भर करती है, न \omterm{lem:b3:forms:partition}{इकाई विभाजन} पर। इसके अतिरिक्त, यदि
$\Phi$ $\Phi(M) = M'$ के साथ दो अभिविन्यस्त उपबहुविधों के प्रतिवेशों की कोई
अवकल समरूपता हो जो $M$ के प्रत्येक घटक के किसी बिंदु पर $M$
की किसी प्रत्यक्ष चौखट को $M'$ की प्रत्यक्ष चौखट पर ले जाती हो,
तो $\int_{M'}\omega =
\int_M\Phi^*\omega$।
\end{lemma}

\begin{proof}
(क) मान लीजिए $\gamma\colon V \to M$, $\delta\colon V' \to M$ ऐसे प्रत्यक्ष प्राचलन हैं जिनके
प्रतिबिंबों में $\operatorname{supp}\omega\cap M$ है। संक्रमण $\tau =
\delta^{-1}\circ\gamma$ $V, V'$ के संगत विवृत
उपसमुच्चयों के बीच अवकल समरूपता है (संक्रमणों की चिकनाई: स्थानीय
आलेख वर्णन के माध्यम से \cref{thm:b3:submanifolds:characterizations}), और वहाँ $\gamma = \delta\circ\tau$, अतः
$\gamma^*\omega = \tau^*(\delta^*\omega)$
(\cref{thm:b3:forms:pullback}(क))। $\delta^*\omega =
g\,\dd u_1\wedge\dots\wedge\dd u_k$ लिखिए; तब (\cref{thm:b3:forms:pullback}(ग)) $\tau^*(\delta^*\omega) =
(g\circ\tau)\,\det(D\tau)\,\dd u_1\wedge\dots\wedge\dd u_k$। दोनों चौखटें
प्रत्यक्ष होने के कारण $D\tau$ किसी धनात्मक आधार को धनात्मक आधार पर
भेजता है: अतः $\det D\tau > 0$, इसलिए $\det D\tau =
\abs{\det D\tau}$, और चर परिवर्तन प्रमेय (\cref{thm:b3:product:changeofvar})
$\int(g\circ\tau)\abs{\det D\tau} = \int g$ दे देती है: अर्थात् दोनों समाकल सहमत हैं। \emph{\omterm{def:b3:forms:orientation}{अभिविन्यास}
के अस्तित्व का पूरा कारण यही है}: चिह्न नियंत्रण के बिना याकोबीय
और उसका निरपेक्ष मान भिन्न हो जाते हैं और समाकल अपरिभाषित हो जाता
है।
(ख) यदि $(\chi_i)$ और $(\tilde\chi_j)$ दो स्वीकार्य विभाजन (आवरण सहित) हों, तो
(क) और परिमित योज्यता से $\sum_i\int\chi_i\omega =
\sum_{i,j}\int\chi_i\tilde\chi_j\omega =
\sum_j\int\tilde\chi_j\omega$, और प्रत्येक द्विक पद किसी भी
चार्ट में परिकलनीय है। अंतिम कथन: यदि $\gamma$ $M$ के प्रत्यक्ष
प्राचलनों में विचरे, तो $\Phi\circ\gamma$ $M'$ के प्रत्यक्ष प्राचलनों में
विचरता है (\omterm{def:b3:forms:orientation}{अभिविन्यास} तुलना स्थानीय रूप से अचर है, और प्रति घटक एक
बिंदु पर नियत की गई है), और $(\Phi\circ\gamma)^*\omega =
\gamma^*(\Phi^*\omega)$।
\end{proof}

\begin{lemma}[इकाई विभाजन, संहत स्थिति]
\label{lem:b3:forms:partition}
मान लीजिए $K \subseteq \R^n$ \omterm{def:b3:topology:compact}{संहत} है और \omterm{def:b3:topology:topology}{विवृत समुच्चय} $W_1, \dots, W_m$ $K$ को आच्छादित
करते हैं। तब $0 \leq \chi_i \leq 1$, \omterm{def:b3:topology:compact}{संहत} $\operatorname{supp}\chi_i \subseteq W_i$, और $K$ के किसी \omterm{def:b3:topology:topology}{प्रतिवेश} पर
$\sum\chi_i = 1$ वाले $\chi_1, \dots, \chi_m \in
\mathcal C^\infty(\R^n)$ विद्यमान हैं।\index{इकाई का विभाजन}
\end{lemma}

\begin{proof}
प्रत्येक $x \in K$ किसी $W_{i(x)}$ में होता है जिसके साथ कोई \omterm{def:b3:forms:closedexact}{संवृत} गोला
$\bar B(x, 2r_x) \subseteq W_{i(x)}$; और संहतता ऐसे $x_1, \dots, x_N$ निकाल देती है कि गोले $B(x_s, r_{x_s})$ $K$ को
आच्छादित कर दें। प्रत्येक $s$ के लिए कोई उभार $\theta_s \in \mathcal
C^\infty$ लीजिए
जिसमें $0 \leq \theta_s \leq 1$, $\bar B(x_s, r_{x_s})$ पर $\theta_s = 1$, और $\operatorname{supp}\theta_s
\subseteq B(x_s, 2r_{x_s})$ (त्रिज्या $\frac32r_{x_s}$ के गोले
के सूचक का मृदुकरण कीजिए, \cref{thm:b3:lp:regularization})। प्रत्येक $s$ को $B(x_s, 2r_{x_s}) \subseteq W_{i(s)}$ वाले
किसी एक सूचकांक $i(s)$ को सौंपिए और $\Theta_i = \sum_{i(s) = i}\theta_s$ रखिए। \omterm{def:b3:topology:topology}{विवृत समुच्चय}
$\Omega_0 = \{\sum_j\Theta_j > \tfrac12\} \supseteq K$ पर फलन $\Theta_i/\sum_j\Theta_j$ काम कर देते हैं, पर वे केवल वहीं परिभाषित हैं;
वैश्विक बनाने के लिए $\rho \in \mathcal
C^\infty(\R^n)$ ऐसा लीजिए कि जहाँ $\sum_j\Theta_j
\geq 1$ वहाँ $\rho = 0$ और
जहाँ $\sum_j\Theta_j \leq \tfrac12$ वहाँ $\rho > 0$ ($1 - \sum_j\Theta_j$ के किसी उपयुक्त कर्तन का मृदुकरण
कीजिए), और
\[
\chi_i = \frac{\Theta_i}{\rho + \sum_j\Theta_j} .
\]
रखिए। हर जगह हर $> 0$ है और $\{\sum_j\Theta_j \geq 1\}$ पर वह $\sum_j\Theta_j$ के बराबर है, जो
$K$ का एक विवृत \omterm{def:b3:topology:topology}{प्रतिवेश} है ($K$ का प्रत्येक बिंदु किसी गोले
$B(x_s, r_{x_s})$ में होता है जहाँ $\theta_s = 1$); और वहाँ $\sum_i\chi_i = 1$। आलंब और परिबंध
स्पष्ट हैं।
\end{proof}

\section{स्टोक्स की प्रमेय}

\begin{definition}\label{def:b3:forms:boundary}
\emph{परिसीमा सहित $k$-उपबहुविध}\index{परिसीमा सहित उपबहुविध}
$M \subseteq \R^n$ वह समुच्चय है जो दो प्रकार के नियमित प्राचलनों से आच्छादित
हो: \emph{आंतरिक चार्ट} $\gamma\colon V \to M \cap W$ जिनमें $V \subseteq \R^k$ विवृत है, और
\emph{परिसीमा चार्ट} $\gamma\colon V \cap H^k \to M
\cap W$, जहाँ $H^k = \{u \in \R^k : u_k \geq 0\}$ और $\gamma$ विवृत $V$ तक
चिकनाई और नियमितता के साथ बढ़ जाता है। \emph{परिसीमा} $\partial M$ उन
बिंदुओं का समुच्चय है जो $u_k = 0$ पर पहुँचे जाते हैं; वह परिसीमा
रहित $(k-1)$-उपबहुविध है, जो प्रतिचित्रणों $u' \mapsto \gamma(u',
0)$ से प्राचलित होती
है। $M$ का कोई \omterm{def:b3:forms:orientation}{अभिविन्यास} \emph{बहिर्मुखी-अभिलंब-पहले} नियम से
$\partial M$ पर एक \omterm{def:b3:forms:orientation}{अभिविन्यास} प्रेरित करता है: $p \in \partial M$ पर $T_p\partial M$ का कोई
आधार $(w_1, \dots, w_{k-1})$ धनात्मक है तभी और केवल तभी जब $(\nu, w_1, \dots, w_{k-1})$ $T_pM$ का धनात्मक
आधार हो, जहाँ $\nu \in T_pM \setminus T_p\partial M$ $M$ से बाहर की ओर संकेत करता है (किसी
परिसीमा चार्ट में: $\nu =
-\partial_k\gamma$, स्पर्शी घटक जोड़ने तक --- \omterm{def:b3:forms:orientation}{अभिविन्यास}
वर्ग उन्हें नहीं देखता)।
\end{definition}

\begin{omfigure}
\begin{tikzpicture}[scale=1.05]
  \fill[omDef!10] (-2,0) arc (180:0:2) -- cycle;
  \draw[thick, omDef] (-2,0) arc (180:0:2);
  \draw[thick, omThm] (-2,0) -- (2,0);
  \node[omDef] at (0.1,1.35) {$M$};
  \node[omThm, below] at (1.55,-0.05) {$\partial M$};
  \draw[->, thick, omExo] (1.415,1.415) -- (2.05,2.05);
  \node[omExo, right] at (2.02,2.12) {$\nu$};
  \draw[->, thick] (1.29,1.53) arc (50:75:2);
  \draw[->, thick, omExo] (-0.6,0) -- (-0.6,-0.65);
  \node[omExo, below] at (-0.6,-0.65) {$\nu$};
  \draw[->, thick] (-0.35,0) -- (0.55,0);
  \node[below] at (0.1,-0.14) {प्रेरित};
\end{tikzpicture}

{\small बहिर्मुखी-अभिलंब-पहले नियम: प्रत्येक \omterm{def:b3:forms:boundary}{परिसीमा} बिंदु पर
बहिर्मुखी सदिश $\nu$ को पहले रखिए; जो आधार उसे $M$ की धनात्मक
चौखट तक पूरा करते हैं वही $\partial M$ को अभिविन्यस्त करते हैं। मानक
\omterm{def:b3:forms:orientation}{अभिविन्यास} वाले किसी समतलीय प्रांत के लिए यह ग्रीन--रीमान का
वामावर्त नियम ही है।}
\end{omfigure}

\begin{lemma}[अर्ध-समष्टि पर स्टोक्स]
\label{lem:b3:forms:halfspace}
मान लीजिए $\eta$ $\R^k$ पर \omterm{def:b3:topology:compact}{संहत} आलंब वाला कोई चिकना $(k-1)$-रूप है।
तब
\[
\int_{H^k}\dd\eta = \int_{\partial H^k}\eta,
\]
जहाँ $H^k$ $\R^k$ का मानक \omterm{def:b3:forms:orientation}{अभिविन्यास} धारण करता है और $\partial H^k = \{u_k = 0\} \cong \R^{k-1}$ प्रेरित
\omterm{def:b3:forms:orientation}{अभिविन्यास}, जो $\R^{k-1}$ के मानक \omterm{def:b3:forms:orientation}{अभिविन्यास} का $(-1)^k$ गुना है।
\end{lemma}

\begin{proof}
पहले \omterm{def:b3:forms:orientation}{अभिविन्यास} का लेखा-जोखा: $\partial H^k$ के अनुदिश बहिर्मुखी अभिलंब
$-e_k$ है, और
\[
\det(-e_k, e_1, \dots, e_{k-1}) =
-\det(e_k, e_1, \dots, e_{k-1}) = -(-1)^{k-1}\det(e_1,
\dots, e_k) = (-1)^k :
\]
अर्थात् $\partial H^k$ की चौखट $(e_1, \dots, e_{k-1})$ प्रेरित \omterm{def:b3:forms:orientation}{अभिविन्यास} के लिए ठीक तब
धनात्मक है जब $k$ सम हो, जिससे कहा गया मिलान निकलता है।
रैखिकता से $\eta
= f\,\dd u_1\wedge\dots\wedge\widehat{\dd
u_i}\wedge\dots\wedge\dd u_k$, $f \in \mathcal
C^\infty_c(\R^k)$ लीजिए; तब $\dd\eta =
(-1)^{i-1}\,\partial_if\,\dd u_1\wedge\dots\wedge\dd u_k$ ($\dd u_i$ को उसके स्थान तक
ले जाने पर $i - 1$ अदला-बदलियाँ लगती हैं)। दो स्थितियाँ, और दोनों
टोनेली--फुबिनी (\cref{thm:b3:product:fubini}) तथा एक-चर वाली कलन की मूल प्रमेय से।

\emph{स्थिति $i < k$।} पहले $\R$ पर $u_i$ में समाकलन करने पर:
$\int_\R\partial_if\,\dd u_i = 0$ (\omterm{def:b3:topology:compact}{संहत} आलंब), अतः $\int_{H^k}\dd\eta = 0$। और $\{u_k = 0\}$ पर $\eta$ के प्रतिबंध में
गुणनखंड $\dd u_k$ होता है, जो $0$ तक प्रतिबंधित हो जाता है (वहाँ
$u_k$ अचर है): अतः $\int_{\partial H^k}\eta =
0$ भी।

\emph{स्थिति $i = k$।} पहले $\intco0\infty$ पर $u_k$ में समाकलन करने पर:
\begin{align*}
\int_{H^k}\dd\eta &= (-1)^{k-1}\int_{\R^{k-1}}
\Bigl(\int_0^\infty\partial_kf\,\dd u_k\Bigr)\dd u' \\
&= (-1)^{k-1}\int_{\R^{k-1}}\bigl(0 - f(u', 0)\bigr)\dd u'
= (-1)^k\int_{\R^{k-1}}f(u',0)\,\dd u' .
\end{align*}
\omterm{def:b3:forms:boundary}{परिसीमा} वाली ओर $\eta$ $f(u',
0)\,\dd u_1\wedge\dots\wedge\dd u_{k-1}$ तक प्रतिबंधित हो जाता है, और प्रेरित
\omterm{def:b3:forms:orientation}{अभिविन्यास} मानक \omterm{def:b3:forms:orientation}{अभिविन्यास} का $(-1)^k$ गुना होने के कारण $\int_{\partial H^k}\eta =
(-1)^k\int_{\R^{k-1}}f(u',0)\,\dd u'$।
दोनों पक्ष सहमत हैं।
\end{proof}

\begin{theorem}[स्टोक्स]\label{thm:b3:forms:stokes}
मान लीजिए $M \subseteq \R^n$ \omterm{def:b3:forms:boundary}{परिसीमा} सहित कोई \omterm{def:b3:topology:compact}{संहत} अभिविन्यस्त
$k$-उपबहुविध है, $\partial M$ प्रेरित \omterm{def:b3:forms:orientation}{अभिविन्यास} धारण करती है, और
$\omega$ $M$ के किसी \omterm{def:b3:topology:topology}{प्रतिवेश} पर कोई चिकना $(k-1)$-रूप। तब\index{स्टोक्स
की प्रमेय}
\[
\int_M\dd\omega = \int_{\partial M}\omega .
\]
विशेष रूप से, यदि $\partial M = \varnothing$ हो: $\int_M\dd\omega = 0$।
\end{theorem}

\begin{proof}
\omterm{def:b3:topology:compact}{संहत} $M$ को परिमित संख्या के प्रत्यक्ष चार्टों (आंतरिक या
\omterm{def:b3:forms:boundary}{परिसीमा}) के प्रतिबिंबों से आच्छादित कीजिए, और $\R^n$ के संगत विवृत
समुच्चयों $W_i$ के अधीनस्थ \omterm{lem:b3:forms:partition}{इकाई विभाजन} $(\chi_i)$ लीजिए (\cref{lem:b3:forms:partition}),
जिसमें $M$ के किसी \omterm{def:b3:topology:topology}{प्रतिवेश} पर $\sum\chi_i
= 1$ हो। उस \omterm{def:b3:topology:topology}{प्रतिवेश} पर $\dd(\sum\chi_i) = 0$,
अतः
\[
\int_M\dd\omega = \sum_i\int_M\dd(\chi_i\omega),
\qquad
\int_{\partial M}\omega =
\sum_i\int_{\partial M}\chi_i\omega :
\]
और दोनों पक्ष योज्य हैं, इसलिए केवल किसी एक ही चार्ट प्रतिबिंब में
आलंबित रूप के लिए प्रमेय सिद्ध करना पर्याप्त है।

\emph{आंतरिक चार्ट।} यदि $\R^k$ में विवृत $V$ के साथ $\operatorname{supp}\omega \cap M
\subseteq \gamma(V)$ हो:
$\eta =
\gamma^*\omega$ को $\R^k$ तक शून्य से बढ़ाइए (चिकना, $V$ में \omterm{def:b3:topology:compact}{संहत} आलंब) और
$\partial H^k$ से दूर आलंब के साथ \cref{lem:b3:forms:halfspace} लगाइए ($V$ को विवृत ऊपरी
अर्ध-समष्टि में स्थानांतरित कीजिए --- अथवा केवल स्थिति $i < k$ का
परिकलन पूरे $\R^k$ पर दोहराइए): अतः $\int_M\dd\omega =
\int_{\R^k}\dd\eta = 0$, और $\int_{\partial M}\omega = 0$ क्योंकि $\omega$
$\partial M$ के निकट लुप्त हो जाता है।

\emph{\omterm{def:b3:forms:boundary}{परिसीमा चार्ट}।} यदि $\operatorname{supp}\omega\cap M
\subseteq \gamma(V \cap H^k)$ हो: शून्य से बढ़ाए गए $\eta = \gamma^*\omega$ के
साथ
$\gamma^*(\dd\omega) = \dd\eta$
(\cref{thm:b3:forms:pullback}(ख)), अतः \cref{def:b3:forms:integral} से:
\[
\int_M\dd\omega = \int_{H^k}\dd\eta
\overset{\text{\cref{lem:b3:forms:halfspace}}}{=}
\int_{\partial H^k}\eta .
\]
अब दाईं ओर को $\int_{\partial
M}\omega$ से पहचानना शेष है। \omterm{def:b3:forms:boundary}{परिसीमा} $\beta(u') =
\gamma(u', 0)$ से प्राचलित
है, और $\beta^*\omega$ $\{u_k = 0\}$ पर $\eta$ का प्रतिबंध है (अंतर्वेशन $u'
\mapsto (u', 0)$ के
साथ $\gamma$ के संयोजन के अंतर्गत \omterm{def:b3:forms:pullback}{प्रत्यानयन})। \omterm{def:b3:forms:orientation}{अभिविन्यास} तुलना दोनों
ओर वही $(-1)^k$ है: चौखट $(\partial_1\beta, \dots, \partial_{k-1}\beta)$ $\partial M$ के प्रेरित \omterm{def:b3:forms:orientation}{अभिविन्यास} में चार्ट
में परिकलित चिह्न $\det(-e_k, e_1, \dots, e_{k-1}) = (-1)^k$ के साथ बैठती है (बहिर्मुखी सदिश का
\omterm{def:b3:forms:pullback}{प्रत्यानयन} $-e_k$ है), जो ठीक वही चिह्न है जो $\partial H^k$ के प्रेरित
\omterm{def:b3:forms:orientation}{अभिविन्यास} को मानक $\R^{k-1}$ से जोड़ता है (\cref{lem:b3:forms:halfspace})। दोनों चिह्न
परिपाटियाँ निरस्त हो जाती हैं: अतः $\int_{\partial H^k}\eta = \int_{\partial M}\omega$।
\end{proof}

\begin{example}[शास्त्रीय प्रमेय]
\label{ex:b3:forms:classical}
मान लीजिए $D \subseteq \R^2$ चिकनी \omterm{def:b3:forms:boundary}{परिसीमा} वाला कोई \omterm{def:b3:topology:compact}{संहत} प्रांत है, जो मानक
रूप से अभिविन्यस्त है। $\omega = P\,\dd x +
Q\,\dd y$ के लिए: $\dd\omega = \bigl(\partial_xQ -
\partial_yP\bigr)\dd x\wedge\dd y$, और स्टोक्स इस रूप में
पढ़ी जाती है
\[
\int_D\Bigl(\frac{\partial Q}{\partial x} -
\frac{\partial P}{\partial y}\Bigr)\dd x\,\dd y
= \oint_{\partial D}P\,\dd x + Q\,\dd y :
\]
अर्थात् \emph{ग्रीन--रीमान}, जो दूसरे~वर्ष के खंड में प्रारंभिक
प्रांतों के लिए सिद्ध हुई थी और अब स्वाभाविक व्यापकता में है।
$\R^3$ में किसी \omterm{def:b3:topology:compact}{संहत} प्रांत पर किसी सदिश क्षेत्र के अभिवाह $2$-रूप
पर लगाई गई स्टोक्स \emph{अपसरण प्रमेय} $\int_\Omega\operatorname{div}F =
\int_{\partial\Omega}\langle F, \nu\rangle\,\dd S$ देती है, और \omterm{def:b3:forms:boundary}{परिसीमा}
सहित किसी पृष्ठ पर किसी $1$-रूप पर लगाई गई शास्त्रीय
\emph{केल्विन--स्टोक्स} कर्ल प्रमेय; \cref{exo:b3:forms:7} दोनों शब्दकोश विस्तार
से लिखता है।
\end{example}

\section{घूर्णन संख्या}\label{sec:b3:forms:winding}

\begin{definition}\label{def:b3:forms:winding}
मान लीजिए $\gamma\colon\intcc01\to\R^2\setminus\{a\}$ कोई चिकना \omterm{def:b3:forms:closedexact}{संवृत} वक्र है। $a$ के परितः उसकी
\emph{घूर्णन संख्या}\index{घूर्णन संख्या} है
\[
\operatorname{Ind}_\gamma(a) =
\frac1{2\pi}\int_\gamma\omega_\theta^a,
\qquad
\omega_\theta^a = \frac{(x - a_1)\,\dd y - (y - a_2)\,\dd
x}{(x - a_1)^2 + (y - a_2)^2} .
\]
\end{definition}

\begin{proposition}\label{prop:b3:forms:windinginteger}
$\operatorname{Ind}_\gamma(a) \in \Z$; और $a$ के फलन के रूप में वह $\R^2\setminus\gamma(\intcc01)$ के प्रत्येक \omterm{def:b3:topology:components}{संबद्ध घटक}
पर अचर है तथा अपरिबद्ध घटक पर शून्य। $\gamma(t) = a + r(\cos2\pi t, \sin2\pi t)$ के लिए: $\operatorname{Ind}_\gamma(a) = 1$।
\end{proposition}

\begin{proof}
$a = 0$ लीजिए और $\gamma = (x, y)$, $\rho =
\norm\gamma > 0$, $\vartheta(t) =
\int_0^t\gamma^*\omega_\theta$ लिखिए, जिससे $\vartheta' =
\frac{xy' - yx'}{x^2 + y^2}$।
सम्मिश्र संकेतन में $u(t)
= \gamma(t)\,\rho(t)^{-1}\eu^{-\iu\vartheta(t)}$ लीजिए; तब $\abs u = 1$, और सीधा परिकलन
\[
\frac{u'}{u} = \frac{\gamma'}{\gamma} - \frac{\rho'}{\rho}
- \iu\vartheta'
= \frac{\bar\gamma\gamma'}{\abs\gamma^2} -
\frac{\rho'}{\rho} - \iu\,\frac{xy' -
yx'}{\rho^2}
= \frac{(xx' + yy')}{\rho^2} - \frac{\rho'}{\rho} = 0,
\]
देता है, क्योंकि $\bar\gamma\gamma' = (xx' + yy') + \iu(xy' - yx')$ और $\rho\rho' = xx' + yy'$। अतः $u$ अचर है: $\abs c = 1$ के साथ
$\gamma(t) =
\rho(t)\,c\,\eu^{\iu\vartheta(t)}$, और $\rho(1) = \rho(0)$ के साथ $\gamma(1) = \gamma(0)$ से $\eu^{\iu\vartheta(1)} = \eu^{\iu\vartheta(0)} = 1$ अनिवार्य हो जाता है:
इसलिए $\vartheta(1) \in 2\pi\Z$, अर्थात् $\operatorname{Ind}_\gamma(0) \in \Z$। \omterm{def:b3:topology:compact}{संहत} वक्र के विवृत पूरक पर $a$ के
फलन के रूप में परिभाषक समाकल \omterm{def:b3:topology:continuity}{संतत} है (\cref{thm:b3:lebesgue:paramcont}, प्रत्येक $a$ के
\omterm{def:b3:topology:topology}{प्रतिवेश} पर प्रभुत्व); और कोई \omterm{def:b3:topology:continuity}{संतत} पूर्णांक-मान वाला फलन स्थानीय
रूप से अचर होता है, अतः घटकों पर अचर। बड़े $\norm a$ के लिए समाकल्य
$O(1/\norm a)$ में एकसमान रूप से $t$ है, अतः सूचकांक $0$ की ओर जाता है
और अपरिबद्ध घटक पर लुप्त हो जाता है। वृत्त के लिए: सीधे $\gamma^*\omega_\theta^a = 2\pi\,\dd t$।
\end{proof}

\begin{remark}
$\C \cong \R^2$ के अंतर्गत $\frac{\dd z}{z} = \frac{x\,\dd x +
y\,\dd y}{x^2 + y^2} + \iu\,\omega_\theta$, अतः $\operatorname{Ind}_\gamma(a) =
\frac1{2\iu\pi}\oint_\gamma\frac{\dd z}{z - a}$: यही \cref{ch:b3:residues} का सूचकांक है, और
\cref{prop:b3:forms:windinginteger} उसकी पूर्णांकता तथा स्थानीय अचरता को वास्तविक-चर साधनों से
पुनः सिद्ध कर देती है --- अर्थात् \omterm{def:b3:residues:singularities}{अवशेष} प्रमेय का संस्थितिक आधा
भाग, जो अब स्टोक्स पर खड़ा है।
\end{remark}

\begin{proposition}\label{prop:b3:forms:exactloop}
यदि $\omega = \dd f$ विवृत $U$ पर \omterm{def:b3:forms:closedexact}{यथार्थ} हो और $\gamma\colon\intcc01\to U$ कोई \omterm{def:b3:forms:closedexact}{संवृत} वक्र हो, तो
$\int_\gamma\omega = 0$।
\end{proposition}

\begin{proof}
$\int_\gamma\dd f = \int_0^1(f\circ\gamma)'(t)\,\dd t =
f(\gamma(1)) - f(\gamma(0)) = 0$ --- शृंखला नियम $\gamma^*(\dd f)$ की पहचान $(f\circ\gamma)'\,\dd t$ से कर देता है।
\end{proof}

\begin{method}\label{met:b3:forms:method}
रूपों के साथ परिकलन: (1) यंत्रवत चलिए --- \omterm{def:b3:forms:wedge}{बाह्य गुणन} चिह्नों के
साथ पुनःक्रमित होते हैं, $\dd$ गुणांकों का अवकलन करके नए $\dd
x_j$
निकालता है, और \omterm{def:b3:forms:pullback}{प्रत्यानयन} प्रतिस्थापन कर देते हैं; बीजगणित पर
भरोसा कीजिए, वह प्रत्येक याकोबीय को कूट रूप में रखता है। (2) किसी
उपबहुविध पर किसी रूप के समाकलन के लिए: प्रत्यक्ष रूप से प्राचलन
कीजिए, \omterm{def:b3:forms:pullback}{प्रत्यानयन} लीजिए, और गुणांक का समाकलन कीजिए; एकमात्र जाल
\omterm{def:b3:forms:orientation}{अभिविन्यास} है --- कोई एक चौखट जाँच लीजिए। (3) किसी समाकल सर्वसमिका
को सिद्ध करने के लिए स्टोक्स की आकृति ढूँढ़िए: क्या समाकल्य
\omterm{def:b3:forms:closedexact}{यथार्थ} है? क्या प्रांत कोई \omterm{def:b3:forms:boundary}{परिसीमा} है? (4) दो ``समांतर''
उपबहुविधों पर समाकलों की तुलना करने के लिए उनके बीच के क्षेत्र पर
स्टोक्स लगाइए (विरूपण तर्क, \cref{exo:b3:forms:10})। (5) किसी \omterm{def:b3:forms:closedexact}{संवृत रूप} का
शून्येतर समाकल किसी संस्थितिक बाधा का प्रमाणपत्र है --- कोई
आद्यंतर नहीं, कोई प्रत्याकर्षण नहीं, कोई शून्य-रहित विस्तार नहीं:
सप्ताहांत समस्या गोले के प्रत्याकर्षणों को इसी प्रकार मारती है।
\end{method}

\section{अभ्यास}

\begin{exercise}[$\star$]\label{exo:b3:forms:1}
$\R^3$ पर मान लीजिए $\omega = x\,\dd y - z\,\dd x$ और $\eta =
\dd x + y\,\dd z$। $\omega\wedge\eta$, $\dd\omega$, $\dd\eta$ और $\dd(\omega\wedge\eta)$
परिकलित कीजिए, और इसी उदाहरण पर श्रेणीबद्ध लाइब्नित्स नियम
सत्यापित कीजिए।
\end{exercise}

\begin{exercise}[$\star$]\label{exo:b3:forms:2}
$\R^3$ पर $\dd$ के तीनों अवतार पहचानिए: $f \in \Omega^0$ के लिए $\dd f \leftrightarrow \nabla f$; कार्य
रूप $\omega_F = F_1\dd x + F_2\dd y + F_3\dd z$ के लिए $\dd\omega_F \leftrightarrow \operatorname{curl}F$; और अभिवाह रूप $\sigma_F = F_1\,\dd y\wedge\dd z + F_2\,\dd
z\wedge\dd x + F_3\,\dd x\wedge\dd y$ के लिए $\dd\sigma_F
\leftrightarrow \operatorname{div}F$। $\dd^2 =
0$
से सर्वसमिकाएँ $\operatorname{curl}\nabla f = 0$ और $\operatorname{div}\operatorname{curl}F = 0$ निकालिए।
\end{exercise}

\begin{exercise}[$\star\star$]\label{exo:b3:forms:3}
तय कीजिए कि प्रत्येक $1$-रूप अपने प्रांत पर \omterm{def:b3:forms:closedexact}{संवृत} है या नहीं,
\omterm{def:b3:forms:closedexact}{यथार्थ} है या नहीं, और जब कोई आद्यंतर हो तब उसे परिकलित कीजिए:
(क) $\R^3$ पर $(2xy + z^2)\,\dd x + x^2\,\dd y + 2xz\,\dd z$;
(ख) $\R^2\setminus\{0\}$ पर $\dfrac{x\,\dd x + y\,\dd y}{x^2 + y^2}$;
(ग) अर्ध-समतल $\{x > 0\}$ पर $\dfrac{-y\,\dd x + x\,\dd y}{x^2 + y^2}$।
\end{exercise}

\begin{exercise}[$\star\star$]\label{exo:b3:forms:4}
(कोणीय रूप) सत्यापित कीजिए कि $\omega_\theta$ (\cref{ex:b3:forms:angular}) \omterm{def:b3:forms:closedexact}{संवृत} है; $0$ के
परितः त्रिज्या $r$ के वृत्त $\gamma$ के लिए $\int_\gamma\omega_\theta$ परिकलित कीजिए;
निष्कर्ष निकालिए कि $\omega_\theta$ $\R^2\setminus\{0\}$ पर \omterm{def:b3:forms:closedexact}{यथार्थ} नहीं है, और यह कि $\R^2\setminus\{0\}$
अपने किसी भी बिंदु के परितः तारकाकार नहीं है (दो मार्ग: \cref{thm:b3:forms:poincare} के
माध्यम से, और सीधे ज्यामिति से)।
\end{exercise}

\begin{exercise}[$\star\star$]\label{exo:b3:forms:5}
(किसी अधिपृष्ठ का क्षेत्रफल रूप) मान लीजिए $M \subseteq \R^n$ कोई \omterm{def:b3:topology:compact}{संहत}
अभिविन्यस्त अधिपृष्ठ है जिसका \omterm{def:b3:forms:orientation}{अभिविन्यास} किसी इकाई अभिलंब क्षेत्र
$\nu$ से दिया गया है ($(w_1, \dots, w_{n-1})$ धनात्मक तभी और केवल तभी जब $\R^n$ में
$(\nu, w_1, \dots, w_{n-1})$ धनात्मक हो)। दर्शाइए कि $(n-1)$-रूप $\sigma_\nu =
\sum_i(-1)^{i-1}\nu_i\,\dd x_1\wedge\dots\wedge\widehat{\dd
x_i}\wedge\dots\wedge\dd x_n$ $M$ पर
\emph{क्षेत्रफल रूप} तक प्रतिबंधित हो जाता है: किसी प्रत्यक्ष
प्राचलन $\gamma$ के लिए $\gamma^*\sigma_\nu = \sqrt{\det G}\,\dd u_1\wedge\dots
\wedge\dd u_{n-1}$, जहाँ $G = ({}^t D\gamma)(D\gamma)$ ग्राम आव्यूह है।
\emph{(ध्यान दीजिए कि $(\gamma^*\sigma_\nu)(e_1, \dots, e_{n-1}) = \det(\nu,
\partial_1\gamma, \dots, \partial_{n-1}\gamma)$, और इस सारणिक का वर्ग कीजिए।)} $\int_{S^2}\sigma_\nu = 4\pi$
परिकलित कीजिए।
\end{exercise}

\begin{exercise}[$\star\star$]\label{exo:b3:forms:6}
$\omega = x\,\dd y\wedge\dd z
+ y\,\dd z\wedge\dd x + z\,\dd x\wedge\dd y$ के लिए $\int_{S^2}\omega$ परिकलित कीजिए: (क) गोलीय निर्देशांकों में सीधे;
(ख) इकाई गोले पर स्टोक्स से। $S^2$ के क्षेत्रफल से $\operatorname{vol}(B^3) = \frac{4\pi}3$ निकालिए,
और सामान्यीकरण कीजिए: $n\operatorname{vol}(B^n) = \operatorname{area}(S^{n-1})$, जो \cref{thm:b3:product:ballvolume} के अनुरूप है।
\end{exercise}

\begin{exercise}[$\star\star$]\label{exo:b3:forms:7}
(शब्दकोश) \cref{thm:b3:forms:stokes} से सावधानी से निकालिए: (क) $\R^3$ में अपसरण प्रमेय
(\cref{exo:b3:forms:2} और \cref{exo:b3:forms:5} को मिलाइए); (ख) $\R^3$ में \omterm{def:b3:forms:boundary}{परिसीमा} सहित किसी \omterm{def:b3:topology:compact}{संहत}
अभिविन्यस्त पृष्ठ के लिए केल्विन--स्टोक्स प्रमेय $\int_S\langle\operatorname{curl}F, \nu\rangle\,\dd S =
\oint_{\partial S}\langle F, \tau\rangle\,\dd\ell$। जाँचिए कि
विषुवत रेखा से घिरे ऊपरी अर्ध-गोलक पर \omterm{def:b3:forms:orientation}{अभिविन्यास} परिपाटियाँ सहमत
हैं।
\end{exercise}

\begin{exercise}[$\star\star$]\label{exo:b3:forms:8}
(ग्रीन सर्वसमिकाएँ) चिकनी \omterm{def:b3:forms:boundary}{परिसीमा} वाले किसी \omterm{def:b3:topology:compact}{संहत} प्रांत $\Omega \subseteq \R^n$ के
\omterm{def:b3:topology:topology}{प्रतिवेश} पर चिकने $u, v$ के लिए सिद्ध कीजिए
\[
\int_\Omega\bigl(u\,\Delta v + \langle\nabla u, \nabla
v\rangle\bigr) = \int_{\partial\Omega}u\,\partial_\nu
v\,\dd S,
\qquad
\int_\Omega\bigl(u\,\Delta v - v\,\Delta u\bigr) =
\int_{\partial\Omega}\bigl(u\,\partial_\nu v -
v\,\partial_\nu u\bigr)\dd S .
\]
निष्कर्ष निकालिए: $\Omega$ पर प्रसंवादी और $\partial\Omega$ पर लुप्त होने वाला
कोई फलन सर्वथा लुप्त हो जाता है, और समान \omterm{def:b3:forms:boundary}{परिसीमा} मानों वाले दो
\omterm{prop:b3:conformal:harmonicholo}{प्रसंवादी फलन} एक ही होते हैं --- अर्थात् \cref{ch:b3:conformal} की दिरिक्ले समस्या
में अद्वितीयता।
\end{exercise}

\begin{exercise}[$\star\star$]\label{exo:b3:forms:9}
(चर परिवर्तन, अभिविन्यस्त रूप) मान लीजिए $\varphi\colon U
\to V$ $\det
D\varphi > 0$ वाली $\R^n$
के विवृत समुच्चयों की कोई अवकल समरूपता है, और $f$ $V$ में \omterm{def:b3:topology:compact}{संहत}
आलंब के साथ \omterm{def:b3:topology:continuity}{संतत}। दर्शाइए कि \omterm{def:b3:forms:pullback}{प्रत्यानयन} सर्वसमिका $\int_U\varphi^*(f\,\dd x_1\wedge\dots\wedge\dd x_n) =
\int_Vf\,\dd x_1\wedge\dots\wedge\dd x_n$ ऐसे $\varphi$
के लिए चर परिवर्तन प्रमेय (\cref{thm:b3:product:changeofvar}) के तुल्य है, और ठीक-ठीक समझाइए
कि याकोबीय पर का निरपेक्ष मान कहाँ चला गया।
\end{exercise}

\begin{exercise}[$\star\star\star$]\label{exo:b3:forms:10}
(विरूपण) मान लीजिए $\omega$ $\R^3\setminus\{0\}$ पर कोई \emph{\omterm{def:b3:forms:closedexact}{संवृत}} $2$-रूप है
और $S_r$ $0$ पर केंद्रित त्रिज्या $r$ का गोलक। दर्शाइए कि
$\int_{S_r}\omega$ $r > 0$ पर निर्भर नहीं करता \emph{(दो त्रिज्याओं के बीच के
खोल पर स्टोक्स लगाइए; दोनों \omterm{def:b3:forms:boundary}{परिसीमा} अभिविन्यासों का ध्यान रखिए)}।
इसे \emph{घन-कोण रूप}
\[
\omega = \frac{x\,\dd y\wedge\dd z + y\,\dd z\wedge\dd x +
z\,\dd x\wedge\dd y}{(x^2 + y^2 + z^2)^{3/2}} :
\]
पर लगाइए; जाँचिए कि वह \omterm{def:b3:forms:closedexact}{संवृत} है, $\int_{S_r}\omega = 4\pi$ परिकलित कीजिए, और निष्कर्ष
निकालिए कि वह $\R^3\setminus\{0\}$ पर \omterm{def:b3:forms:closedexact}{संवृत} है पर \omterm{def:b3:forms:closedexact}{यथार्थ} नहीं --- $\omega_\theta$ का
द्वि-विमीय सहोदर, और स्थिरवैद्युतिकी में गाउस के नियम की ज्यामितीय
अंतर्वस्तु।
\end{exercise}

\begin{exercise}[$\star\star$]\label{exo:b3:forms:11}
मान लीजिए $\gamma$ $\R^2\setminus\{0\}$ में $n = \operatorname{Ind}_\gamma(0)$ वाला कोई चिकना \omterm{def:b3:forms:closedexact}{संवृत} वक्र है।
दर्शाइए कि $\R^2\setminus\{0\}$ पर \emph{प्रत्येक} \omterm{def:b3:forms:closedexact}{संवृत} $1$-रूप $\omega$ के लिए
$\int_\gamma\omega = n\int_{S^1}\omega$ \emph{(\cref{exo:b3:forms:12} से $\omega = c\,\omega_\theta
+ \dd f$ लिखिए)}। व्याख्या: छिद्रित समतल पर
आवर्तों के लुप्त होने की एकमात्र बाधा \omterm{def:b3:forms:winding}{घूर्णन संख्या} है।
\end{exercise}

\begin{exercise}[$\star\star\star$]\label{exo:b3:forms:12}
(पहली दे राम गणना) दर्शाइए कि $U = \R^2\setminus\{0\}$ पर प्रत्येक \omterm{def:b3:forms:closedexact}{संवृत} $1$-रूप
$\omega$ अद्वितीय रूप से
\[
\omega = c\,\omega_\theta + \dd f,
\qquad c = \frac1{2\pi}\int_{S^1}\omega,\quad f \in
\mathcal C^\infty(U) :
\]
है: $(1,0)$ से $p$ तक किसी \omterm{def:b3:topology:pathconnected}{पथ} (पहले त्रिज्य टुकड़ा, फिर वृत्तीय
चाप) के अनुदिश $\omega - c\,\omega_\theta$ का समाकलन करके $f(p)$ परिभाषित कीजिए, दर्शाइए
कि परिणाम चयनों से स्वतंत्र है और वह भी ठीक इसलिए कि $S^1$-आवर्त
लुप्त हो जाता है, और $\dd
f = \omega - c\,\omega_\theta$ जाँचिए। निष्कर्ष निकालिए: $H^1(\R^2\setminus\{0\}) \cong \R$, जो
कोणीय रूप से जनित है।
\end{exercise}

\section{समस्या: ब्राउवर की स्थिर-बिंदु प्रमेय}

\begin{problem}[{सप्ताहांत समस्या --- कोई प्रत्याकर्षण नहीं, कोई
पलायन नहीं}]\label{pb:b3:forms:1}
ब्राउवर की प्रमेय कहती है कि \emph{\omterm{def:b3:forms:closedexact}{संवृत} इकाई गोले $\bar B = \bar B^n \subseteq \R^n$ का उसी
में जाने वाला प्रत्येक \omterm{def:b3:topology:continuity}{संतत प्रतिचित्रण} किसी स्थिर बिंदु को धारण
करता है} --- गणित की महान प्रमेयों में से एक, जिसके परिणाम खेल
सिद्धांत (नैश साम्यावस्थाएँ) से लेकर आव्यूह विश्लेषण तक फैले हैं।
अवकल-रूप उपपत्ति सबसे स्वच्छ ज्ञात उपपत्ति है: स्टोक्स दिखा देती
है कि गोलक गोले का प्रत्याकर्षण नहीं है, और शेष सब उसी से निकलता
है। सर्वत्र $S = S^{n-1} = \partial\bar B$, $n \geq 2$, और
\[
\sigma = \sum_{i=1}^n(-1)^{i-1}x_i\,
\dd x_1\wedge\dots\wedge\widehat{\dd x_i}\wedge\dots\wedge
\dd x_n
\]
(टोपी गुणनखंड को हटा देती है)।

\medskip
\noindent\textbf{भाग I --- मापक यंत्र।}
\begin{enumerate}
  \item $\dd\sigma$ परिकलित कीजिए, और स्टोक्स (\cref{thm:b3:forms:stokes}) से $\int_S\sigma =
        n\operatorname{vol}(\bar B) > 0$ निष्कर्ष
        निकालिए, जहाँ $S$ गोले की \omterm{def:b3:forms:boundary}{परिसीमा} \omterm{def:b3:forms:orientation}{अभिविन्यास} धारण करता
        है।
  \item $n = 2$ और $n = 3$ के लिए $S$ पर $\sigma$ के प्रतिबंध की पहचान
        चापलंबाई और क्षेत्रफल रूपों से कीजिए ($\nu(x) = x$ के साथ \cref{exo:b3:forms:5})
        और $\int_S\sigma$ सीधे पुनः परिकलित कीजिए।
  \item मान लीजिए $W \subseteq \R^N$ विवृत है और $\varphi\colon W
        \to \R^n$ ऐसा चिकना प्रतिचित्रण
        कि सभी $x \in W$ के लिए $\norm{\varphi(x)} = 1$। दर्शाइए कि $\varphi^*(\dd\sigma) = 0$।
        \emph{($\norm\varphi^2 = 1$ का अवकलन कीजिए: $D\varphi(x)$ का प्रतिबिंब अधिसमतल
        $\varphi(x)^\perp$ में होता है, जिसकी विमा $n - 1$ है; अब \cref{prop:b3:forms:linearpullback}(2)
        लगाइए।)}
  \item निम्नलिखित ``उपपत्ति'' में कि $\int_S\sigma = 0$ --- ``$S$ \omterm{def:b3:forms:boundary}{परिसीमा}
        रहित \omterm{def:b3:topology:compact}{संहत} है, और $S$ पर प्रतिबंधित $\sigma$ उस पर एक
        उच्चतम-घात रूप है, अतः \omterm{def:b3:forms:closedexact}{संवृत}, अतः स्टोक्स से $\int_S\sigma =
        \int_S\dd(\text{किसी वस्तु}) = 0$'' ---
        दोष कहाँ है? (गलत शब्द पर ठीक-ठीक उँगली रखिए।)
  \item एक अनुच्छेद में भाग II की रणनीति समझाइए: किसका समाकलन
        होगा, किस पर, और विरोधाभास कहाँ से आएगा।
\end{enumerate}

\medskip
\noindent\textbf{भाग II --- कोई चिकना प्रत्याकर्षण नहीं।}
विरोधाभास के लिए मान लीजिए $r$ गोले का उसके गोलक पर एक
\emph{चिकना प्रत्याकर्षण} है: अर्थात् $r$ $\bar B$ के किसी \omterm{def:b3:topology:topology}{प्रतिवेश}
पर चिकना है, $r(\bar B) \subseteq S$, और सभी $x \in S$ के लिए $r(x) = x$।
\begin{enumerate}[resume]
  \item $\int_Sr^*\sigma = \int_S\sigma$ उचित ठहराइए \emph{($S$ पर $r$ तत्समक तक
        प्रतिबंधित हो जाता है: यदि $\gamma$ $S$ के किसी टुकड़े का
        प्रत्यक्ष प्राचलन हो, तो $r\circ\gamma = \gamma$)}।
  \item $\bar B$ पर स्टोक्स का उपयोग करके $\int_Sr^*\sigma =
        \int_{\bar B}\dd(r^*\sigma)$ दर्शाइए।
  \item $\dd(r^*\sigma) = r^*(\dd\sigma) = 0$ दर्शाइए ($\varphi = r$ पर लगाया गया प्रश्न 3), और निष्कर्ष
        निकालिए: \emph{कोई चिकना प्रत्याकर्षण $\bar B \to S$ विद्यमान नहीं
        है।}
  \item अपवर्जित स्थिति $n = 1$ हाथ से निपटाइए: सीधे दर्शाइए कि कोई
        भी \omterm{def:b3:topology:continuity}{संतत प्रतिचित्रण} $\intcc{-1}1 \to
        \{-1, 1\}$ दोनों सिरों को स्थिर नहीं रखता,
        और आपने जिस प्रमेय का उपयोग किया उसका नाम बताइए।
\end{enumerate}

\medskip
\noindent\textbf{भाग III --- चिकना ब्राउवर।} मान लीजिए $g$
$\bar B$ के किसी \omterm{def:b3:topology:topology}{प्रतिवेश} पर चिकना है, $g(\bar B)
\subseteq \bar B$, और $\bar B$ में उसका
\emph{कोई} स्थिर बिंदु नहीं है।
\begin{enumerate}[resume]
  \item $\delta = \min_{x\in\bar B}\norm{g(x) - x} >
        0$ दर्शाइए।
  \item $x \in \bar B$ के लिए मान लीजिए $u(x) = \frac{x -
        g(x)}{\norm{x - g(x)}}$, और मान लीजिए $r(x) = x +
        t(x)\,u(x)$ किरण
        $\{x + tu(x) : t \geq 0\}$ का $S$ से प्रतिच्छेदन है। द्विघात $\norm{x + tu}^2 = 1$ हल कीजिए
        और
        \[
        t(x) = -\langle x, u(x)\rangle +
        \sqrt{1 - \norm x^2 + \langle x, u(x)\rangle^2}
        \;\geq\; 0 .
        \]
        प्राप्त कीजिए।
  \item दर्शाइए कि $\bar B$ पर मूल के भीतर की राशि \emph{कड़ाई से}
        धनात्मक है: यदि $1 - \norm x^2 + \langle x,
        u(x)\rangle^2 = 0$ हो, तो $\norm x = 1$ और $\langle
        x, u(x)\rangle = 0$, अर्थात्
        $\langle x, x -
        g(x)\rangle = 0$, यानी $\langle x, g(x)\rangle =
        1$; और $\norm x = 1$, $\norm{g(x)} \leq 1$ के साथ कोशी--श्वार्ज़
        से इससे $g(x) = x$ अनिवार्य हो जाता है --- जो अपवर्जित है।
        निष्कर्ष निकालिए कि $r$ $\bar B$ के किसी \omterm{def:b3:topology:topology}{प्रतिवेश} पर चिकना
        है।
  \item $r(\bar B) \subseteq S$ दर्शाइए और $x
        \in S$ के लिए $r(x) = x$ भी \emph{($\norm x = 1$ के
        लिए $\langle x, u(x)\rangle \geq 0$ का उपयोग करते हुए $t(x) = 0$ जाँचिए, जो स्वयं $\langle x, x - g(x)\rangle = 1 -
        \langle x, g(x)\rangle \geq 0$
        से निकलता है)}। भाग II के साथ निष्कर्ष निकालिए:
        \emph{$\bar B$ के प्रत्येक चिकने स्व-प्रतिचित्रण का कोई
        स्थिर बिंदु होता है।}
\end{enumerate}

\medskip
\noindent\textbf{भाग IV --- \omterm{def:b3:topology:continuity}{संतत} ब्राउवर।} मान लीजिए $f\colon\bar B\to\bar B$ स्थिर
बिंदु रहित \omterm{def:b3:topology:continuity}{संतत प्रतिचित्रण} है।
\begin{enumerate}[resume]
  \item $\varepsilon = \min_{\bar B}\norm{f - x} > 0$ दर्शाइए, और $\sup_{\bar B}\norm{p - f} <
        \varepsilon/2$ वाला कोई \emph{बहुपदीय}
        प्रतिचित्रण $p\colon\R^n\to
        \R^n$ बनाइए (स्टोन--वाइरश्ट्रास, \cref{thm:b3:complete:stoneweierstrass},
        निर्देशांक-दर-निर्देशांक --- अदिश से सदिश आसन्नन तक जाना
        उचित ठहराइए)।
  \item प्रतिचित्रण $p$ गोले से बाहर जा सकता है; अतः $g = \frac{p}{1
        + \varepsilon/2}$
        रखिए। $g(\bar B) \subseteq \bar B$ और $\sup_{\bar B}\norm{g - f} < \varepsilon$ दर्शाइए।
  \item भाग III के साथ विरोधाभास निकालिए और निष्कर्ष निकालिए:
        \emph{प्रत्येक \omterm{def:b3:topology:continuity}{संतत प्रतिचित्रण} $\bar B^n \to \bar B^n$ का कोई स्थिर बिंदु
        होता है।}
  \item उदाहरण से दर्शाइए कि यह प्रमेय इन पर विफल हो जाती है:
        विवृत गोला; गोलक $S$; कोई \omterm{def:b3:forms:closedexact}{संवृत} वलय। प्रत्येक
        प्रतिउदाहरण $\bar B$ का कौन-सा गुण खो देता है?
\end{enumerate}

\medskip
\noindent\textbf{भाग V --- लाभांश।}
\begin{enumerate}[resume]
  \item (पेरों--फ्रोबेनियस, अस्तित्व) मान लीजिए $A$ ऐसा $n\times n$
        आव्यूह है जिसकी सभी प्रविष्टियाँ $> 0$ हैं, और $\Delta = \{x \in \R^n : x_i \geq 0,\ \sum x_i =
        1\}$।
        दर्शाइए कि प्रतिचित्रण $x \mapsto Ax/\norm{Ax}_1$ $\Delta$ पर सुपरिभाषित और \omterm{def:b3:topology:continuity}{संतत}
        है, कि $\Delta$ $\R^{n-1}$ के किसी \omterm{def:b3:forms:closedexact}{संवृत} गोले का समस्थितिक है
        (अपने आफ़ीन विस्तार में अरिक्त \omterm{def:b3:topology:interior}{अंतःभाग} वाले किसी उत्तल
        \omterm{def:b3:topology:compact}{संहत} से त्रिज्य \omterm{def:b3:topology:continuity}{समस्थितिकता}), और निष्कर्ष निकालिए कि
        $A$ का कोई ऐसा अभिलक्षणिक सदिश है जिसकी सभी प्रविष्टियाँ
        कड़ाई से धनात्मक हैं और जिसका अभिलक्षणिक मान $> 0$ है।
  \item निष्कर्ष निकालिए कि धनात्मक प्रविष्टियों वाले प्रत्येक
        प्रसंभाव्य आव्यूह (जिसके स्तंभों का योग $1$ हो) का कोई
        स्थायी प्रायिकता सदिश $\pi = A\pi$ होता है --- अर्थात्
        पेजरैंक-प्रकार का सदिश। (अद्वितीयता भी लागू होती है, पर
        उसके लिए अन्य साधन चाहिए।)
  \item (रोमिल गोलक, व्यवस्था) मान लीजिए $v$ $S = S^{n-1}$ के किसी
        \omterm{def:b3:topology:topology}{प्रतिवेश} पर चिकना है, जिसमें $x \in S$ के लिए $\langle v(x),
        x\rangle = 0$ और $\norm{v(x)} = 1$
        (एक \emph{इकाई स्पर्श क्षेत्र})। $t \in \R$ के लिए $F_t(x) = x + t\,v(x)$
        रखिए। $S$ पर $\norm{F_t(x)} = \sqrt{1 + t^2}$ दर्शाइए: अर्थात् $F_t$ $S$ को
        गोलक $\sqrt{1+t^2}\,S$ में भेजता है।
  \item दर्शाइए कि $P(t) = \int_SF_t^*\sigma$ $t$ में एक \emph{बहुपद} है \emph{(किसी
        चार्ट में $F_t^*\sigma$ का प्रत्येक गुणांक फलन $t$ में बहुपदीय
        है, जिसके गुणांक चार्ट चर में चिकने हैं; और समाकलन रैखिक
        है)}।
  \item दर्शाइए कि छोटे $\abs t$ के लिए $F_t$ $S$ से $\sqrt{1+t^2}\,S$ पर एक
        अवकल समरूपता है: $t\operatorname{Lip}(v) < 1$ के लिए एकैकीपन; प्रतिलोम फलन
        प्रमेय से स्थानीय अवकल समरूपता (चार्टों में \cref{thm:b3:submanifolds:ift}); और
        प्रतिबिंब \omterm{def:b3:topology:connected}{संबद्ध} लक्ष्य गोलक में विवृत तथा \omterm{def:b3:forms:closedexact}{संवृत}। \cref{lem:b3:forms:consistency}
        और $x \mapsto \lambda x$ के अंतर्गत मापन $\sigma_{\lambda x} =
        \lambda^{n}\,\sigma_x$ (उसे जाँचिए) का उपयोग करते
        हुए निष्कर्ष निकालिए कि
        \[
        P(t) = \pm(1 + t^2)^{n/2}\int_S\sigma,
        \qquad\text{चिह्न } + \text{ के साथ, छोटे } t \text{ के लिए}
        \]
        (\omterm{def:b3:forms:orientation}{अभिविन्यास} $t = 0$ से \omterm{def:b3:topology:continuity}{सांतत्य} द्वारा सुरक्षित)।
  \item निष्कर्ष निकालिए (मिल्नोर): यदि $n$ \emph{विषम} हो, तो
        $(1 +
        t^2)^{n/2}$ $t$ में बहुपद नहीं है, फिर भी वह $0$ के निकट
        बहुपद $P(t)/\int_S\sigma$ से सहमत है --- विरोधाभास। अतः \emph{सम-विमीय
        गोलक $S^{n-1}$ ($n$ विषम) कोई इकाई स्पर्श क्षेत्र धारण नहीं
        करते}, और मानकीकरण तथा मृदुकरण से (घटकशः संवलन और प्रक्षेप
        --- दोनों चरण उचित ठहराइए) कोई \omterm{def:b3:topology:continuity}{संतत} कहीं-भी-अलुप्त स्पर्श
        क्षेत्र भी नहीं: अर्थात् पृथ्वी पर बहने वाली हर हवा कोई न
        कोई शांत बिंदु छोड़ जाती है।

  \item (विषम गोलक स्वतंत्रता से सँवरते हैं) $S^{2m-1}
        \subseteq \R^{2m} \cong \C^m$ पर एक स्पष्ट
        चिकना इकाई स्पर्श क्षेत्र प्रस्तुत कीजिए: सम्मिश्र संकेतन
        में $v(x) = \iu x$, अर्थात्
        \[
        v(x_1, y_1, \dots, x_m, y_m) =
        (-y_1, x_1, \dots, -y_m, x_m) .
        \]
        स्पर्शिता और इकाई लंबाई सत्यापित कीजिए, और निष्कर्ष
        निकालिए कि प्रश्न 23 का सम-विषमता द्विभाजन तीक्ष्ण है:
        कोई गोलक ठीक तभी सँवारा जा सकता है जब उसकी विमा विषम हो।
        सम $n$ के लिए प्रश्न 22 का बहुपदीय तर्क कहाँ टूट जाता
        है?
  \item (\omterm{def:b3:forms:boundary}{परिसीमा} व्यवहार से शून्य) मान लीजिए $f \colon \bar
        B^n \to \R^n$ \omterm{def:b3:topology:continuity}{संतत} है और
        प्रत्येक $x \in S^{n-1}$ के लिए $\langle f(x),
        x\rangle \geq 0$। दर्शाइए कि $f$ $\bar B^n$ में
        कहीं लुप्त हो जाता है। \emph{(यदि नहीं, तो $g(x) = -f(x)/\norm{f(x)}$ $\bar B$ को
        संततता से $S \subseteq \bar B$ में भेजता है; $g$ पर ब्राउवर लगाइए और
        \omterm{def:b3:forms:boundary}{परिसीमा} परिकल्पना का खंडन कीजिए।)} \emph{आच्छादकता कसौटी}
        निकालिए: $\norm x \to \infty$ पर $\frac{\langle F(x), x\rangle}{\norm x} \to
        +\infty$ वाला कोई \omterm{def:b3:topology:continuity}{संतत} $F\colon\R^n\to\R^n$ आच्छादक होता
        है --- अरैखिक विश्लेषण के प्रबलकारिता तर्कों का
        परिमित-विमीय पूर्वज।
\end{enumerate}
\end{problem}
